\scanpage{20}

une forme quadratique à valeurs dans le module inversible $\Delta'\inserted{=\Delta^\vee}$. Il revient au $=$ de se donner

\[
\Delta\longrightarrow\underline{\operatorname{Quad}}(E)\qquad
\left(\simeq\underline{\operatorname{End}}(E)/\underline O\cdot1\unsure{)}\deleted{\unsure{\otimes}}\otimes L^\vee\right)
\]

(vrai pour \unsure{tt} $\Delta$ loc. libre de r. f.)

i.e. de se donner

\[
u_Q:\underbrace{\deleted{\unsure{M}}\Delta\otimes L}\longrightarrow\underline{\operatorname{End}}(E)/\underline O\cdot1
\]

Nous allons en profiter pour définir une alg. \deleted{\unsure{pr.}} quadratique $\mathcal A_Q$ dans $X$, et un hom \inserted{d'alg.}

\[
\mathcal A_Q\mathbin{\unsure{-|\longrightarrow}}\underline{\operatorname{End}}(E)
\]

dont l'image soit l'anneau $\mathcal A_\Sigma$ des similitudes pour la structure de similitude $\Sigma\inserted{=\Sigma_Q}$ définie par $Q$, \inserted{i.e l'image inv. de $\operatorname{Im}u_Q$ dans $\underline{\operatorname{End}}(E)$.} Pour ceci, on \deleted{\unsure{regarde}} la structure d'extension

\[
0\longrightarrow\underline O_X\longrightarrow\underline{\operatorname{End}}(E)\longrightarrow\underline{\operatorname{End}}(E)/\underline O\cdot1_E\longrightarrow0
\]

et \inserted{on désigne par $\mathcal A_Q$} l'image inverse de cette extension par $u_Q$, donc on a un hom d'extensions

\diagram{\begin{tikzcd}[ampersand replacement=\&,column sep=small]
0\arrow[r,no head]\&\underline O\arrow[r,no head]\arrow[d,"{\mathrm{id}}"]\&\mathcal A_Q\arrow[r,no head]\arrow[d,"{u'_Q}"]\&\overset{\Delta\otimes L}{\overset{\scriptscriptstyle\Vert\,\mathrm{déf\unsure{n}}}{\Lambda}}\arrow[r,no head]\arrow[d,"{u_Q}"]\&0\\
0\arrow[r,no head]\&\underline O\arrow[r]\&\underline{\operatorname{End}}(E)\arrow[r,no head]\&\operatorname{End}(E)/\underline O\arrow[r]\&0
\end{tikzcd}}
% archive sheet number: 19
